\documentclass[border=6pt]{standalone}
% NOTE: the `tikz` class option is deliberately NOT used: with it, standalone
% falls back to a full letter page for circuitikz pictures (Overfull \hbox at
% \end{circuitikz}) instead of cropping to the drawing.  Plain border=6pt crops
% tightly around the drawing, as in the target.
\usepackage[european]{circuitikz}
\usetikzlibrary{arrows.meta}   % solid triangular head for the current arrow
% The target uses american symbols: a zigzag resistor and a four-bump coil,
% even though the package is loaded with the european option.
% A bump of the american coil is (bipoles/americaninductor/coil height) *
% (bipoles/length) tall: 0.15*1.4cm = 0.21cm by default, which reads flat next
% to the 0.42cm R_f zigzag.  Raising it to .21 gives ~0.29cm, matching the
% target's four-loop coil, which is comparable in height to the resistor.
\ctikzset{resistor=american, inductor=american}
\ctikzset{bipoles/americaninductor/coil height=0.21}

\begin{document}
\begin{circuitikz}

  % ==================================================================
  % Geometry (all in cm).  Drawing spans x = 0 .. 12.1 and y = -0.95 .. 3.3,
  % i.e. 12.1 x 4.25 (aspect 2.85; rendered page aspect ~2.68, target ~2.70).
  %
  %   lower (return) wire  y = 0
  %   upper (loop) wire    y = 2.2      -> the two wires are 2.2 cm apart
  %   top dashed border    y = 3.3      -> 1.1 cm headroom above the upper wire
  %   bottom dashed border y = -0.95    -> 0.95 cm below the return wire
  %
  % Horizontal distribution as fractions of the 12.1 cm width
  % (target values in brackets):
  %   converter block  0.041-0.248  [0.043-0.255]
  %   L_f              0.322-0.430  [0.324-0.433]
  %   R_f              0.467-0.545  [0.468-0.548]
  %   current arrow    0.574-0.620  [0.577-0.623]
  %   port marks +/v/- 0.628        [0.629]
  %   group divider    0.690-0.698  [~0.700]  (crosses bare wire only)
  %   Z_g              0.769-0.855  [0.770-0.860] (wholly inside the grid box)
  %   source / right wire 0.921     [0.924]
  % All component labels share one row at y = 2.85 (0.65 cm above the wire),
  % which leaves ~0.2 cm to the top border and ~0.1 cm to the coil tops.
  % ==================================================================

  % --- converter AC-side terminals (on the block's right edge) --------
  \coordinate (cTR) at (3.0,2.2);    % upper wire leaves the right edge
  \coordinate (cBR) at (3.0,0);      % lower wire leaves the bottom-right corner

  % --- filter and grid element ends -----------------------------------
  \coordinate (l1) at (3.9,2.2);     % L_f left end
  \coordinate (l2) at (5.2,2.2);     % L_f right end
  \coordinate (r1) at (5.65,2.2);    % R_f left end
  \coordinate (r2) at (6.6,2.2);     % R_f right end
  \coordinate (z1) at (9.3,2.2);     % Z_g left end  (inside the grid box)
  \coordinate (z2) at (10.35,2.2);   % Z_g right end (inside the grid box)
  \coordinate (gN) at (11.15,2.2);   % grid side, upper corner
  \coordinate (vT) at (11.15,1.5);   % v_g upper terminal
  \coordinate (vB) at (11.15,0.4);   % v_g lower terminal
  \coordinate (gS) at (11.15,0);     % grid side, lower corner

  % --- series loop: converter -> L_f -> R_f -> Z_g -> v_g -> converter -
  % Long bare run from the port marks to Z_g, short run from Z_g to the
  % source, exactly as in the target.  Component labels are placed
  % separately below so that they all sit on one row.
  \draw (cTR) -- (l1) to[L] (l2)
               -- (r1) to[R] (r2)
               -- (z1) to[generic] (z2)
               -- (gN) -- (vT)
               to[sV] (vB)
               -- (gS) -- (cBR);

  % --- DC/AC converter block (plain TikZ: no circuitikz component) ----
  % Square 2.5 x 2.5, ~45 degree diagonal; bottom edge on the return wire,
  % top edge 0.3 cm above the upper wire.
  \draw (0.5,0) rectangle (3.0,2.5);
  \draw (0.5,0) -- (3.0,2.5);        % DC-to-AC diagonal, bottom-left -> top-right
  \node at (1.2,1.93) {DC};          % upper-left side of the diagonal
  \node at (2.35,0.35) {AC};         % lower-right side, near the bottom-right corner

  % --- dashed group boxes with their captions ------------------------
  % 0.1 cm gap between the GFM and grid boxes at the divider x = 8.35-8.45;
  % the divider crosses bare wire only (Z_g starts at x = 9.3).
  \draw[dashed] (0.0,-0.95) rectangle (8.35,3.3);   % GFM converter group
  \draw[dashed] (8.45,-0.95) rectangle (12.1,3.3);  % grid group
  \node at (5.0,-0.6)    {GFM converter};           % 0.41 of the width (target 0.42)
  \node at (10.275,-0.6) {grid};                    % 0.85 of the width (target 0.85)

  % --- component labels: one common row, clear of border and divider --
  \node at (4.55,2.85)  {$L_f$};     % over the coil
  \node at (6.125,2.85) {$R_f$};     % over the zigzag
  \node at (9.825,2.85) {$Z_g$};     % over the grid impedance, same row

  % --- annotations ---------------------------------------------------
  % current arrow on the upper wire, pointing to the grid side
  \draw[-{Triangle[length=6pt,width=5pt]}] (6.95,2.2) -- (7.5,2.2);
  \node at (7.22,2.85) {$i$};
  % port voltage v: v, + and - share one x, well inside the GFM box
  \node at (7.6,1.9) {$+$};          % just under the upper wire
  \node at (7.6,1.1) {$v$};          % midway between the two wires
  \node at (7.6,0.3) {$-$};          % just above the lower wire
  % grid voltage source label, level with the circle centre, to its left
  \node[anchor=east] at (10.72,0.95) {$v_g$};

\end{circuitikz}
\end{document}
